\section{Model dictionary, Round31 checkpoint and Round32 plan}
\label{sec:model}

\subsection{Model dictionary}
\label{sec:dictionary}

The theory is the SU(2) Kogut--Susskind form on the spatial lattice
$\mathbb Z^3$ at fixed spacing $a$. A coarse site $b\in\mathbb Z^3$ owns the
three positive links from each fine tail $(4b_x+r,2b_y+q,b_z)$, $0\le r<4$,
$0\le q<2$: exactly 24 untruncated $L^2(\mathrm{SU}(2))$ link factors. Both
original endpoint gauge actions are retained. Write $C_e=-\sum_aX_{ea}^2$ for
the link Casimir, with generators $i\sigma_a/2$, and $W_f=\tfrac12\tr U_f$ for
the real fundamental Wilson plaquette, so $\norm{W_f}=1$.

An elementary face in directions $a<c$ at $p$ has positive link tails at $p$,
$p+e_a$, $p+e_c$ and $p$, so its coarse owner set is
$\{\pi(p),\pi(p+e_a),\pi(p+e_c)\}$; there is no tail at $p+e_a+e_c$. Grouping
by anchor gives 24 face classes per coarse site
(Table~\ref{tab:dict-classes}), with relative supports inside the star
$S=\{0,e_x,e_y,e_z\}$ (\repo{research/round21/forward/i1/report.md}{I1 dictionary}). Three classes lie inside one factor. They
form the ten-link selected strip, with coefficient triple
$(\lambda_L,\mu_M,\lambda_R)$ constrained by
$|\lambda_L|,|\lambda_R|\le\alpha/2$ and $|\mu_M|\le\alpha/8$; the other 14
links of the factor are free. The remaining 21 classes cross a block boundary.
They are the \emph{omitted} classes, with the common coefficient
$\nu=\alpha\tau/24$.

\begin{table}[htbp]
\centering\small
\caption{The I1 face classes per coarse site; $q$ is the fine $y$ phase.}
\label{tab:dict-classes}
\begin{tabular}{lccl}
\toprule
Orientation and phase & Classes & Relative support & Role\\
\midrule
$xy$: $r=0,1,2$; $q=0$ & 3 & $\{0\}$ & selected\\
$xy$: $r=0,1,2$; $q=1$ & 3 & $\{0,e_y\}$ & omitted\\
$xy$: $r=3$; $q=0$ & 1 & $\{0,e_x\}$ & omitted\\
$xy$: $r=3$; $q=1$ & 1 & $\{0,e_x,e_y\}$ & omitted\\
$xz$: $r=0,1,2$; $q=0,1$ & 6 & $\{0,e_z\}$ & omitted\\
$xz$: $r=3$; $q=0,1$ & 2 & $\{0,e_x,e_z\}$ & omitted\\
$yz$: $r=0,\dots,3$; $q=0$ & 4 & $\{0,e_z\}$ & omitted\\
$yz$: $r=0,\dots,3$; $q=1$ & 4 & $\{0,e_y,e_z\}$ & omitted\\
\bottomrule
\end{tabular}
\end{table}

Throughout Sections~\ref{sec:av1}--\ref{sec:aw2} the selected triple is
exactly $(0,0,0)$. This is an
admissible choice inside the closed AQ coefficient box, not a replacement of a
nonzero selected vacuum by Haar. It is also not the uniform model: equality
with uniform Kogut--Susskind theory would require all 24 coefficients to equal
$\alpha\tau/24$, which fails for $\tau\ne0$. The original AL dictionary
$\tau=96/g^4$ is not assigned to this patterned subfamily; Section~\ref{sec:ax1}
assigns it to the uniform model.

Whole stars $b+S$ are retained only when $b+S\subset\Lambda_N=[-N,N]^3$. In
normalized units $\delta=\alpha/8$ the I1.5 convention groups the omitted faces
by anchor:
\begin{equation}
 \phi_b=-\frac{\tau}{3}\sum_{f\in O_b}W_f,\quad |O_b|=21,\quad
 \norm{\phi_b}=7|\tau|,\qquad
 \widehat H_N=\sum_{b\in\Lambda_N}h_b+\sum_{b+S\subset\Lambda_N}\phi_b,\quad
 h_b=8\!\!\sum_{e\text{ owned by }b}\!\!C_e .
 \label{eq:dict-star}
\end{equation}
At the zero triple the on-site operator is a pure Casimir sum with constant
(Haar) vacuum projection $P_b$, zero ground scalar and $h_b\ge6(I-P_b)$. The
physical Hamiltonian is $H_N=\delta\widehat H_N$. In the energy unit $\alpha$
each star is $V_b=\phi_b/8=-(\tau/24)\sum_{f\in O_b}W_f$, with
$\norm{V_b}\le7|\tau|/8$. Mixing the two unit systems is an eightfold clock
error, and every Round32 checker rejects it.

For the original $xz$ Wilson plaquette $W$ at the origin, the complete-factor
cover is $R=\{0,e_z\}$: 48 stored links and 36 original endpoints (22 in each
factor, eight shared). Its seven incident whole-star anchors are
\begin{equation}
 R-S=\{0,-e_x,-e_y,-e_z,e_z,e_z-e_x,e_z-e_y\}.
 \label{eq:dict-seven}
\end{equation}
All seven are retained for $N\ge2$. The two-anchor orthant count, the four
drawn links of $W$ as a cover, and a one-star interaction budget are rejected
controls.

The on-site product reference on $R$ is the Haar vector $\Omega_R$ with
projection $P_R$. Weyl integration gives
$E_{\rm Haar}[W^n]=1,0,\tfrac14,0,\tfrac18$ for $n=0,\dots,4$, so the reference
moments are $\omega_0(W)=0$ and $\omega_0(W^2)=\tfrac14$. The vector
$W\Omega_R$ carries spin $\tfrac12$ on four links with Casimir $\tfrac34$
each; its energy is $8\cdot4\cdot\tfrac34\,\delta=24\delta=3\alpha$. The free
centered spectral measure is therefore $\tfrac14\delta_3$ in units of
$\alpha$, and
\begin{equation}
 C_0(s)=\tfrac14e^{-3s},\qquad c_0(\theta)=\tfrac14e^{3i\theta}.
 \label{eq:dict-reference}
\end{equation}
These formulas describe the reference, not the interacting spectrum.

\paragraph{AM2 expansion constants.}
AM2 (Round29) proves, for every finite complete-factor volume and both signs of
$|\tau|\le10^{-8}$, a unique ground and a centered gap at least $\alpha/16$
(normalized $\tfrac12$; \repo{research/round29/advisor/am2-gate.json}{AM2 gate}). The proof uses a nilpotent commuting
creation-operator expansion. For nonempty $I\subset\Lambda$ a creation vector
$c_I\in\bigotimes_{x\in I}Q_x\mathcal H_x$, with $Q_x=I-P_x$, defines
$\hat c_I=|c_I\rangle\langle\Omega_I|\otimes1$. These operators commute, and
$\hat c_I\hat c_J=0$ whenever $I\cap J\ne\varnothing$. With
$C=\sum_I\hat c_I$ and the anchored norm
$\norm{c}_a=\max_u\sum_{I\ni u}\norm{c_I}$, the ground vector is
$\psi=e^{-C}\Omega_0$, with vacuum coefficient one, and the unique fixed point
satisfies
\begin{equation}
 c=\sum_{k=0}^{8}\frac{L_k(c,\ldots,c)}{k!},\qquad
 (L_k)_M=H_M^{-1}P_M\,\mathrm{ad}_{C_1}\cdots\mathrm{ad}_{C_k}(V)\,\Omega_0,\qquad
 \norm{c}_a\le R_a=\tfrac1{64}.
 \label{eq:dict-fixed}
\end{equation}
The commutator nest terminates at order eight. The per-site interaction sum
counts all four stars that contain a site, including incoming stars:
$J=\max_u\sum_{X\ni u}\norm{V_X}\le4\cdot7|\tau|=28|\tau|\le
J_0=7/25000000$. The all-order majorant and its evaluated constants are
\begin{equation}
 \begin{gathered}
 G(t)=16e^{8t}(1+10t),\qquad G(R_a)<\tfrac{148}{7},\qquad G'(R_a)<352,\\
 J_0G(R_a)<\tfrac{37}{6250000}<R_a,\qquad 2J_0G'(R_a)<1 .
 \end{gathered}
 \label{eq:dict-majorant}
\end{equation}
The fixed-point equation gives the two majorant inequalities
$\norm{c}_a\le J\,G(\norm{c}_a)$ and
$\norm{c-L_0}_a\le J\bigl(G(\norm{c}_a)-16\bigr)$, which AV1 uses. AM2 removes
the on-site representation cutoff for the eigenvalues and the gap; AV1 removes
it for the ground vector. We write $R_a$ for AM2's anchored radius so that $R$
always denotes the cover.

\paragraph{State and clock.}
AQ1 (Round29) supplies, by diagonal extraction along centered whole-star boxes,
locally normal stationary subsequential states on the full coarse algebra, with
trace-norm convergence of the local densities on every finite region. Let
$\omega_\tau$ be any such state for the declared coefficients, $\Omega_\tau$ its
GNS vacuum and $H_\tau\ge0$ its physical energy generator. AQ2 gives a gap at
least $\alpha/16$ in that representation; the Round32 window certificate uses
only $H_\tau\ge0$. No uniqueness, whole-sequence convergence or rate in $N$ is
assumed. Both the reference and the interacting clocks use
\begin{equation}
 G_\tau=H_\tau/\alpha,\qquad s=\frac{\alpha t_E}{\hbar},\qquad
 \theta=\frac{\alpha t}{\hbar}.
 \label{eq:dict-clock}
\end{equation}
The positive energy $\alpha$, action $\hbar$, spacing $a$ and comparison energy
$E_\star$ stay fixed, so $s$, $\theta$ and $\tau$ are dimensionless. The older
normalized clock $u=s/8$, with free Wilson exponent 24, is forbidden in Round32
evidence packets. For the original plaquette,
\begin{equation}
 m_\tau=\omega_\tau(W),\quad
 \chi_\tau=(\pi_\tau(W)-m_\tau)\Omega_\tau,\quad
 C_\tau(s)=\langle\chi_\tau,e^{-sG_\tau}\chi_\tau\rangle,\quad
 c_\tau(\theta)=\langle\chi_\tau,e^{i\theta G_\tau}\chi_\tau\rangle .
 \label{eq:dict-observable}
\end{equation}
The correlation is dimensionless and centered as a vector in the actual state.
Changing $\tau$ between investigations would be an explicit interaction change,
not a change of units; in sub-rounds~1 and~2 the coupling is not changed.

\subsection{Round31 checkpoint}
\label{sec:checkpoint}

Round31 closed three investigations in the same zero-selected family
(Table~\ref{tab:checkpoint}; \repo{papers/round31-addendum/main.pdf}{Round31
addendum}). AT4 kept the original cap and missed its target. AT5 and AT6 met
the target only after a prospectively declared change of coupling to
$10^{-14}$, which selects a different Hamiltonian and its own AQ state.

\begin{table}[htbp]
\centering\small
\caption{Round31 results inherited as the starting point; decimals summarize
outward rational bounds in the frozen outputs.}
\label{tab:checkpoint}
\begin{tabularx}{\linewidth}{@{}l>{\raggedright\arraybackslash}p{0.23\linewidth}>{\raggedright\arraybackslash}X l@{}}
\toprule
Loop & Design & Result & Verdict\\
\midrule
AT4 & $\tau=10^{-8}$, $s=1$, $L=10^4$, target $10^{-6}$ &
radius $\approx8.41518704386267\times10^{-4}$; interval
$[0.0116052483875797,\,0.0132882857963523]$ & limited\\[3pt]
AT5 & $\tau=10^{-14}$, $s=1$, $L=10^9$ &
datum $\approx0.0124467670919660$, radius
$<8.086109564653\times10^{-7}$ & accepted\\[3pt]
AT6 & $\tau=10^{-14}$, $0\le s\le128$, 4097 nodes &
every node error $<8.491932\times10^{-7}$;
$I\in[0.083174562675757,\,0.083535946641292]$, containing $1/12$ & accepted\\
\bottomrule
\end{tabularx}
\end{table}

The AT4 budget was dominated by the reset bound
$D=2\sqrt{49|\tau|/3}\approx8.0829\times10^{-4}$, which scales as
$|\tau|^{1/2}$. Before Round32 production, the advisor also observed that
optimizing AT4's Poisson certificate over the cutoff $L$ leaves a
dynamics-plus-tail cost of about $1.27\times10^{-6}$ at $s=1$ even if the state
term vanished. That was a planning observation; AV2 later proves it as a
retained failure (Section~\ref{sec:av2-failures}). The cap therefore had two
separate obstacles: a square-root state term and a kernel whose absolute first
moment diverges. Sub-round~1 addresses one each. All three Round31 intervals
contain their free references, so no interaction shift was resolved there
either. The ranked Round31 roadmap goals were cap-level precision, an
interaction-induced effect, a uniform Hamiltonian, state identification and a
continuum trajectory (\repo{research/round31/advisor/roadmap.json}{Round31 roadmap}).

\subsection{Round32 plan}
\label{sec:plan}

A user instruction of 23 September 2026 authorized Round32 as an adaptive
expert-panel cycle with three deliberation loops and five sub-rounds of two
investigations. In loop~1 each lens wrote a memo, source ledger and
recommendation, and the skeptic wrote a prospective triage and control
catalogue. In loop~2 the panel critiqued plan~v1. It corrected the face count
per factor from the bound 84 to the derived 49, set the AV1 target from AV2
feasibility, separated the AV1 and AV2 routes, and adopted a pre-registration
block. Loop~3 was an objection pass. All four sign-offs arrived without a veto
on the plan; the skeptic's veto on freezing AV1 was cleared once the selection
note matched the contract. Plan~v2 was frozen before AV1 production
(\repo{research/round32/advisor/plan.json}{plan record}). Table~\ref{tab:plan} lists it. All ten rows
are gated.

\begin{table}[htbp]
\centering\small
\caption{Frozen Round32 plan (v2). Direction: P = paired forward/reverse,
S = single direction plus independent skeptic replay, T = statement plus
skeptic.}
\label{tab:plan}
\begin{tabularx}{\linewidth}{@{}>{\raggedright\arraybackslash}p{0.15\linewidth}
 >{\raggedright\arraybackslash}X c
 >{\raggedright\arraybackslash}p{0.17\linewidth}@{}}
\toprule
Sub-round, goal & Investigations & Dir. & Status in this edition\\
\midrule
1, cap &
AV1: first-order local state lemma. &
P & gated: accepted within scope\\
 &
AV2: window-kernel certificate for $C(1)$ at $\tau=10^{-8}$. &
P & gated: accepted within scope; reference unresolved\\[4pt]
2, shift &
AW1: Peter--Weyl parity, link-flip antisymmetry, first-order Wilson mean and
second-order remainder $K_2$, with the AW2 coupling rule frozen in advance. &
P & gated: accepted within scope\\
 &
AW2: sign-certified enclosure of $\omega_\tau(W)$ at both signs, coupling
fixed by the AW1 rule. &
S & gated: accepted within scope\\[4pt]
3, uniform &
AX1: route-B re-derivation for uniform Kogut--Susskind SU(2), $J_0$
re-frozen at $29/10^8$, dictionary $\tau=96/g^4$. &
P & gated: accepted within scope\\
 &
AX2: uniform-model datum $C(1)$ at the cap by the window route. &
S & gated: accepted within scope; reference unresolved\\[4pt]
4, state &
AY1: uniform local closeness of all AQ-type subsequential states and the
boundary-independent first-order density. &
P & gated: accepted within scope; local closeness, not uniqueness\\
 &
AY2: two named construction families, topology, explicit constant, and what
remains unproved. &
T & gated: accepted within scope; local closeness, not uniqueness; static, not dynamic\\[4pt]
5, continuum &
AZ1: regulator/coupling trajectory and the one uniform estimate supplied or
failed. &
T & gated: accepted within scope; one uniform estimate at fixed a; nothing uniform in a\\
 &
AZ2: exact first- and second-order coefficients on the Round11 two-plaquette
graph, labelled a different finite model. &
S & gated: accepted within scope; sign-certified on the finite graph only;
static, not dynamic; no transfer\\
\bottomrule
\end{tabularx}
\end{table}

\paragraph{Direction rules.}
Derivation-bearing investigations are paired, with genuinely different forward
and reverse routes disclosed in the contract. Formula-instantiating
certificates use one direction plus an independent skeptic replay written from
the contract alone. State and continuum statements use a statement plus
skeptical review. Same-author passes are never independent review. For paired
investigations the reverse producer receives only \texttt{AGENTS.md}, the
contract and its declared shared premises; the freezing tool enforces that
inventory. After each sub-round the advisor may reorder, replace or retitle the
next pair, and records any change in the next selection note and the plan
history. After each of sub-rounds~1 to~4 the goals were kept
(\repo{research/round32/advisor/panel-update-1.md}{panel updates~1},
\repo{research/round32/advisor/panel-update-2.md}{2},
\repo{research/round32/advisor/panel-update-3.md}{3} and
\repo{research/round32/advisor/panel-update-4.md}{4}).

\paragraph{Pre-registration.}
Every contract freezes and hashes, before any outcome, the model identifier,
the selected triple, the coupling and the signs evaluated, the state
provenance, the observable with its centering convention, the nodes, an
itemized list of error terms and the target. It also fixes the outcome types
(\recid{accepted_within_scope}, \recid{limited}, \recid{insufficient}), the
allowed sub-labels, the claim exclusions and the required control ids.
Post-hoc node selection is forbidden and pilot runs are excluded from every
count. An error term may be marked not applicable only with a stated reason.
Each checker reads its target and reference from the hash-checked contract
snapshot and records its own SHA-256 before the full-size evaluation. The
allowed sub-labels are \recid{reference_unresolved},
\recid{sign_certified_finite_graph}, \recid{sign_certified_below_cap},
\recid{static_not_dynamic} and
\recid{uniform_local_closeness_not_uniqueness}. Four loop-2 vetoes stand: no
AW2 coupling chosen after $K_2$ is known and no sign margin below 2, no uniform
Wilson-mean sign rider on AX2, no coefficients fitted from finite-graph
enclosures, and no radius target of $10^{-8}$ for AV. A fifth veto, on any
priority or independence sentence about AM2, was resolved by the fixed
attribution sentence of the status section. A limited or insufficient verdict is a
retained outcome. Investigation counts are never a fraction of the continuum
problem.
